\pagebreak
\section{Kirchhoff's Rules}
\hrule
\vspace{.5cm}

\begin{eqboxed}{eq:KirchhoffVoltageRule}{Kirchhoff Voltage Rule}
{In a \textit{\textbf{closed}} circuit, if you return to the same location, no matter what path you take, the change in voltage will be 0. Very similar to conservation of energy in a magnetic field.}
\sum\Delta V_n = 0 
\end{eqboxed}

\begin{eqboxed}{eq:KirchhoffCurrentRule}{Kirchhoff Current Rule}
{At any node (location where two or more paths/wires meet), the current in must be equal to the current out (of said node)}
\boxed{\sum I_{in} = \sum I_{out}}, \text{at any node}
\end{eqboxed}

\begin{example}
    \textbf{Single Loop Example\label{ex:SingleLoopExample}}
    \vspace{.2cm}
    \hrule 
    \vspace{.2cm}

    \begin{center}
        \begin{tikzpicture}[american]
            % Draw the main loop
            \draw (0,0) 
                to[battery1, l=$V_b$, invert] (0,6) % Battery on the left
                to[short, i=$I$, -*] (4,6) node[anchor=south west] {a} % Top wire with current arrow
                to[R, l=$R_1$, v=$+$, -*] (4,3) node[anchor=west] {b} % Resistor 1
                to[R, l=$R_2$, v=$+$, -*] (4,0) node[anchor=north west] {c} % Resistor 2
                to[short, i=$I$] (0,0); % Bottom wire with current arrow
        \end{tikzpicture}

        Can be simply solved by using Ohms law and doing \(I = \frac{V_b}{R_1 +R_2}\), or using Kirchhoff\\ 
        \vspace{.25cm}
        Using Kirchhoff's Voltage Rule (\ref{eq:KirchhoffVoltageRule}) we find the drop in current between \(a \rightarrow b \rightarrow c \rightarrow V_b\)
        \[\Delta V_1 = +IR_1, \quad \Delta V_1 = +IR_2, \quad \Delta V_1 = -V_b\]
        \[ IR_1 + IR_2 - V_b = 0 \quad \rightrightarrows \quad I = \frac{V_b}{R_1 +R_2 }\]
    \end{center}
\end{example}

\begin{example}
    \textbf{Two Loop Example\label{ex:TwoLoopExample}}
    \vspace{.2cm}
    \hrule
    \vspace{.2cm}
    \begin{center}
        
            \begin{tikzpicture}[american, scale=1.2]
                % Coordinates
                \coordinate (bottom_left) at (0,0);
                \coordinate (top_left) at (0,4);
                \coordinate (a) at (4,4);
                \coordinate (b) at (4,0);
                \coordinate (c) at (7,0);
                \coordinate (top_right) at (7,4);

                % Left Loop: V1 and R1 (with polarity)
                \draw (bottom_left) to[battery1, l=$V_1$, invert] (top_left)
                    (top_left) -- (a);
                \draw (b) to[R, l=$R_1$, v^>=$ $] (bottom_left); % v label adds the +/-

                % Middle Branch: R2 (with polarity)
                \draw (b) to[R, l=$R_2$, v_<=$ $] (a);

                % Right Branch: R3 (with polarity) and V2
                \draw (a) -- (top_right)
                    (top_right) to[R, l=$R_3$, v^>=$ $] ++(0,-2)
                    to[battery1, l=$V_2$] (c)
                    (c) -- (b);

                % Nodes a, b, c
                \fill (a) circle (1.5pt) node[above] {a};
                \fill (b) circle (1.5pt) node[below] {b};
                \fill (c) circle (1.5pt) node[below] {c};

                % Current arrows (Blue)
                \draw[white, thick, ->, >=stealth] (1.5, 4.3) -- (2.5, 4.3) node[midway, above] {$I_1$};
                \draw[white, thick, ->, >=stealth] (2.5, -0.75) -- (1.5, -0.75) node[midway, below] {$I_1$};
                \draw[white, thick, ->, >=stealth] (3.6, 3.5) -- (3.6, 2.5) node[midway, left] {$I_2$};
                \draw[white, thick, ->, >=stealth] (7.4, 2.5) -- (7.4, 1.5) node[midway, right] {$I_3$};

            \end{tikzpicture}

        Using Kirchhoff's Current Rule(\ref{eq:KirchhoffCurrentRule}) we find that \(I_1 + I_2 +I_3 = 0\). To solve this we need equations with \(I_1, I_2, I_3\)

        Using Kirchhoff's Voltage Rule (\ref{eq:KirchhoffVoltageRule}) on the \textit{outer} loop we get 
        \[I_3R_3 + V_2 + I_3R_1 - V_1\]
        On the inner/left loop 
        \[I_2R_2 + I_1R_1 -V_1\]
    \end{center}

\end{example}